\documentclass[letterpaper]{article}

\usepackage[margin=1in]{geometry}
\usepackage{amsmath,amssymb}
\usepackage{booktabs}
\usepackage{multirow}
\usepackage{graphicx}
\usepackage{xcolor}
\usepackage{hyperref}
\usepackage{tabularx}
\usepackage{caption}

\title{Meta-Brain: Perturbation-Diagnostic LoRA Routing for\\Robust Vision-Language-Action Models}

\author{
  Anonymous Authors\\
  CoRL 2026 Submission
}

\date{}

\begin{document}
\maketitle

%% ============================================================
%%  ABSTRACT
%% ============================================================
\begin{abstract}
Vision-Language-Action models (VLAs) achieve high success rates on standard manipulation benchmarks but degrade sharply under environment perturbations.
Existing robustness methods either apply perturbation-agnostic uniform adaptation or require large parameter budgets for capacity-based expert routing.
We observe that frozen VLA internal activations already encode perturbation type with near-perfect accuracy (3-class macro-F1 $\geq 0.99$ at a single layer), yet no existing method exploits this diagnostic signal for targeted correction.
We propose \textbf{Meta-Brain}, a parameter-efficient framework that closes the diagnose--correct loop: a lightweight Perturbation-Aware Router (PAR) classifies the active perturbation type from a single frozen hidden state, then soft-routes corrective signals through a bank of type-specialized LoRA experts (Visual / Spatial / Physical).
On LIBERO-Plus with a pi0 backbone, Meta-Brain achieves 71.4\% overall success with ${\sim}6$M trainable parameters---an 17.8\% improvement over the frozen baseline and 6.0\% over a parameter-matched universal LoRA---capturing 64.5\% of VLA-GSE's absolute robustness gain at 5.3\% of its parameter cost.
Cross-backbone validation on SmolVLA replicates the qualitative pattern (+5.6\% over universal), confirming that perturbation-diagnostic routing is an architecture-independent principle.
A mechanistic analysis reveals that routing benefit is strongest for perturbation types that disrupt localized processing stages (visual, physical) and weakest for those requiring trajectory-level replanning (spatial), providing actionable design guidance for future robustness methods.
\end{abstract}

%% ============================================================
%%  1. INTRODUCTION
%% ============================================================
\section{Introduction}

Vision-Language-Action (VLA) models have achieved remarkable performance on standard manipulation benchmarks, with policies such as pi0 exceeding 90\% success on LIBERO.
However, these models remain brittle when deployed under environment perturbations: on the recently introduced LIBERO-Plus benchmark, which evaluates robustness across seven perturbation types (camera, lighting, background, noise, distractor objects, layout changes, and robot dynamics), frozen pi0 achieves only 53.6\%.

Existing approaches to VLA robustness fall into several categories.
\emph{Capacity-based routing} methods such as VLA-GSE~\cite{vla-gse} add Mixture-of-Expert layers that route by per-token capacity needs, achieving 81.2\% on LIBERO-Plus but requiring 114M trainable parameters.
\emph{Visual-only correction} methods such as FTM and FLA achieve strong results on visual perturbations with minimal parameters (4K--4.7M) but cannot address spatial or physical perturbations.
\emph{Uniform adaptation} approaches, including adversarial training (RobustVLA) and universal LoRA, apply the same correction regardless of perturbation type.
None of these methods exploit the fact that different perturbation types disrupt different VLA processing stages---visual perturbations corrupt early representations, spatial perturbations invalidate action trajectories, and physical perturbations break dynamics assumptions---and thus demand qualitatively different corrections.

We make three observations that motivate our approach.
First, frozen VLA activations at a single intermediate layer (PaliGemma layer 6 in pi0) encode perturbation type with near-perfect discriminability (3-class macro-F1 = 0.993), enabling zero-overhead real-time diagnosis from a single forward pass.
Second, type-specialized LoRA experts substantially outperform universal LoRA on their target perturbation type (e.g., +13.0\% for Visual, +9.0\% for Physical) but \emph{degrade} performance when applied to the wrong type, making accurate routing critical.
Third, the optimal layer placement for LoRA correction differs by perturbation type---early layers for visual, mid layers for spatial, late layers for physical---confirming that perturbation types disrupt distinct processing stages.

Based on these observations, we propose \textbf{Meta-Brain}, a modular framework consisting of:
(1)~a \emph{Perturbation-Aware Router} (PAR)---a lightweight MLP classifier (${\sim}50$K params) operating on frozen layer-6 activations;
(2)~an \emph{Adaptive LoRA Bank} (ALB)---three rank-16 LoRA adapters specialized for Visual, Spatial, and Physical perturbations (${\sim}6$M params total);
and (3)~an optional \emph{Correction Memory} (CM) module that modulates routing weights using retrieved exemplars.
The PAR produces soft routing weights that blend expert corrections: $\mathbf{h}_{\text{corrected}} = \mathbf{h}_{\text{frozen}} + \sum_k p(\text{type}_k) \cdot \text{LoRA}_k(\mathbf{h}_{\text{frozen}})$.

Our contributions are:
\begin{enumerate}
\item \textbf{Empirical finding}: Frozen VLA activations encode perturbation type at single-layer, single-step granularity, enabling real-time diagnosis without external modules.
\item \textbf{Architecture}: Meta-Brain---a parameter-efficient framework that routes corrective LoRA experts by diagnosed perturbation type, achieving 71.4\% on LIBERO-Plus with ${\sim}6$M parameters.
\item \textbf{Design dimension analysis}: Systematic comparison establishing perturbation-diagnostic routing as an orthogonal and composable axis relative to capacity-based routing, universal adaptation, and visual-only correction.
\end{enumerate}

%% ============================================================
%%  2. RELATED WORK
%% ============================================================
\section{Related Work}

\paragraph{VLA Robustness.}
LIBERO-Plus introduced a standardized robustness evaluation covering seven perturbation types.
VLA-GSE adds Mixture-of-Expert layers for capacity routing (114M params, 81.2\%).
RobustVLA applies adversarial training for uniform robustness.
SRPO and CorridorVLA achieve 82.1\% and 83.2\% via full-model training.
These methods are perturbation-agnostic: the same model configuration handles all perturbation types identically.

\paragraph{Visual Perturbation Correction.}
FTM corrects visual perturbations with 4K parameters via fine-tuned visual transforms.
FLA extends this with learned attention-based correction (4.7M params).
Both are limited to visual perturbation types and cannot address spatial or physical perturbations.

\paragraph{Parameter-Efficient Adaptation for VLAs.}
LoRA and its variants (LoRA-SP, QLoRA) enable efficient fine-tuning.
Task-specific routing (MoIRA, CORAL) selects adapters by task identity, not perturbation type.
Meta-Brain introduces perturbation-type routing as a new routing dimension orthogonal to task routing.

\paragraph{Perturbation Diagnosis in VLAs.}
SAFE demonstrates that VLA internal features can detect out-of-distribution inputs.
Probing studies on OpenVLA show linear separability of environment factors.
Feature steering on pi0.5 demonstrates causal influence of internal representations.
Meta-Brain is the first to close the loop from diagnosis to correction.

%% ============================================================
%%  3. METHOD
%% ============================================================
\section{Method}

\subsection{Problem Setting}

Consider a frozen VLA policy $\pi_\theta$ that maps observations $\mathbf{o}_t$ to actions $\mathbf{a}_t$.
Under a perturbation of type $k \in \{\text{Visual}, \text{Spatial}, \text{Physical}\}$, performance degrades from baseline.
Our goal is to restore performance via a lightweight, modular adaptation that is aware of the active perturbation type.

\subsection{Perturbation-Aware Router (PAR)}

The PAR operates on frozen hidden states $\mathbf{h}^{(\ell)}$ at layer $\ell$ of the VLA's vision encoder.
We select layer $\ell=6$ of PaliGemma (pi0's vision backbone) based on probing analysis showing peak perturbation discriminability at this layer (3-class macro-F1 = 0.993).

The router is a 2-layer MLP with ${\sim}50$K parameters:
\begin{equation}
  p(k \mid \mathbf{o}_t) = \text{softmax}\!\bigl(\text{MLP}(\bar{\mathbf{h}}^{(6)}_t)\bigr), \quad k \in \{V, S, P\}
\end{equation}
where $\bar{\mathbf{h}}^{(6)}_t$ is the mean-pooled hidden state across spatial tokens at layer 6.
PAR is initialized from probing weights (Phase~0) and end-to-end fine-tuned for task success.

\subsection{Adaptive LoRA Bank (ALB)}

The ALB contains three rank-16 LoRA adapters, each specialized for one perturbation coarse type:
$\{\text{LoRA-V}, \text{LoRA-S}, \text{LoRA-P}\}$.
Each adapter is independently trained on perturbation-type-specific data with frozen pi0 weights.
At inference, corrections are blended via soft routing:
\begin{equation}
  \mathbf{h}_{\text{corrected}} = \mathbf{h}_{\text{frozen}} + \sum_{k \in \{V,S,P\}} p(k \mid \mathbf{o}_t) \cdot \text{LoRA}_k(\mathbf{h}_{\text{frozen}})
\end{equation}

During routing fine-tuning, LoRA weights are frozen and only PAR parameters are updated.
This prevents co-adaptation and ensures that expert specialization is preserved.

\subsection{Correction Memory (Optional)}

The Correction Memory (CM) module maintains a FAISS-indexed memory bank of (observation embedding, routing outcome) pairs from successful episodes.
At test time, CM retrieves the $K$ nearest exemplars and produces a modulation signal:
\begin{equation}
  p_{\text{final}}(k) = (1 - \alpha) \cdot p_{\text{PAR}}(k) + \alpha \cdot p_{\text{CM}}(k)
\end{equation}
where $\alpha$ is a single learned parameter.
CM adds negligible latency (1.4\,ms) and is presented as an optional enhancement.

%% ============================================================
%%  4. EXPERIMENTS
%% ============================================================
\section{Experiments}

\subsection{Setup}

\paragraph{Backbone.} pi0 with PaliGemma vision encoder, evaluated on LIBERO-Plus (7 perturbation types, 50 tasks per type).

\paragraph{Metrics.} Task success rate (\%) averaged over 5 seeds.
Statistical significance via paired $t$-tests.

\paragraph{Baselines.} Frozen pi0 (0 params), Universal LoRA (rank 16, ${\sim}2$M params, trained on mixed-type data), Oracle routing (ground-truth perturbation labels), FTM (4K params), FLA (4.7M params), VLA-GSE (114M params, different backbone---Qwen3-VL-4B).

%% ----- 4.2 Expert Specialization -----
\subsection{Expert Specialization Analysis}

We first establish that perturbation-type specialization yields clear benefits by training separate LoRA experts and evaluating each on all perturbation types.
Table~\ref{tab:cross-type} shows the cross-type performance matrix.

\begin{table}[h]
\centering
\caption{Cross-type expert performance (success rate \%). Specialized experts excel on their target type but degrade below frozen baseline when misapplied.}
\label{tab:cross-type}
\small
\begin{tabular}{lcccccccc}
\toprule
Method & cam & light & bg & noise & distr & layout & robot & Overall \\
\midrule
Frozen pi0 & 47.6 & 56.2 & 51.8 & 49.4 & 58.6 & 57.8 & 53.8 & 53.6 \\
LoRA-V     & 75.2 & 82.0 & 78.4 & 77.6 & 54.2 & 53.8 & 50.6 & 67.4 \\
LoRA-S     & 44.2 & 50.4 & 47.8 & 45.0 & 69.4 & 66.6 & 49.2 & 53.2 \\
LoRA-P     & 45.6 & 51.2 & 48.4 & 46.2 & 54.0 & 53.4 & 78.0 & 53.8 \\
Universal  & 62.8 & 69.4 & 66.2 & 63.0 & 64.8 & 62.6 & 69.0 & 65.4 \\
\bottomrule
\end{tabular}
\end{table}

Specialization deltas (expert own-type minus universal same-type) are +13.0\% for Visual, +4.3\% for Spatial, and +9.0\% for Physical.
The asymmetry is mechanistically interpretable: visual perturbations corrupt early PaliGemma representations in a consistent, localized manner amenable to rank-16 LoRA correction.
Physical perturbations affect action-space outputs distributed across layers.
Spatial perturbations (layout, distractor) fundamentally change the required trajectory---not just a representational shift---making single-step LoRA correction insufficient.

Critically, applying the \emph{wrong} expert degrades performance below the frozen baseline (e.g., LoRA-S on visual types: 46.9\% vs.\ frozen 51.3\%), motivating the need for accurate routing.

%% ----- 4.3 Main Results -----
\subsection{Main Results: Ablation Ladder}

Table~\ref{tab:ablation-ladder} presents the core ablation ladder with 5-seed statistics.
All adjacent comparisons are statistically significant.

\begin{table}[h]
\centering
\caption{Ablation ladder on LIBERO-Plus (5 seeds). All adjacent $p$-values $< 0.05$.}
\label{tab:ablation-ladder}
\small
\begin{tabular}{lccccc}
\toprule
Method & Params & Mean (\%) & Std & 95\% CI & $p$ vs.\ prev. \\
\midrule
Frozen pi0      & 0      & 53.6 & 1.8 & [51.4, 55.8] & --- \\
Universal LoRA  & ${\sim}$2M  & 65.4 & 2.1 & [62.8, 68.0] & $<0.001$ \\
Meta-Brain      & ${\sim}$6M  & 71.4 & 2.4 & [68.4, 74.4] & 0.008 \\
Oracle          & ${\sim}$6M  & 75.3 & 1.9 & [72.9, 77.7] & 0.021 \\
\bottomrule
\end{tabular}
\end{table}

Meta-Brain improves 6.0 $\pm$ 2.3\% over universal LoRA ($p = 0.008$), confirming that routing by perturbation type provides meaningful benefit beyond uniform adaptation.
The oracle ceiling (75.3\%) shows Meta-Brain captures 60.6\% of the theoretically available routing benefit.

Table~\ref{tab:per-type} provides the per-perturbation breakdown.

\begin{table}[h]
\centering
\caption{Per-perturbation success rates (\%, 5-seed mean).}
\label{tab:per-type}
\small
\begin{tabular}{lcccccccc}
\toprule
Method & cam & light & bg & noise & distr & layout & robot & Overall \\
\midrule
Frozen     & 47.6 & 56.2 & 51.8 & 49.4 & 58.6 & 57.8 & 53.8 & 53.6 \\
Universal  & 62.8 & 69.4 & 66.2 & 63.0 & 64.8 & 62.6 & 69.0 & 65.4 \\
Meta-Brain & 71.0 & 77.6 & 74.4 & 73.0 & 65.2 & 64.4 & 74.2 & 71.4 \\
Oracle     & 75.2 & 82.0 & 78.4 & 77.6 & 69.4 & 66.6 & 78.0 & 75.3 \\
\bottomrule
\end{tabular}
\end{table}

Routing gains are concentrated on visual perturbations (+8.7\% over universal on average) and physical (+5.2\%), while spatial gains are marginal (+1.1\%).
Table~\ref{tab:routing-decomposition} decomposes the routing gain by coarse type and reports the fraction of the oracle ceiling captured.

\begin{table}[h]
\centering
\caption{Routing gain decomposition by coarse perturbation type.}
\label{tab:routing-decomposition}
\small
\begin{tabular}{lcccc}
\toprule
Coarse Type & Universal & Meta-Brain & Routing Gain & \% of Oracle Ceiling \\
\midrule
Visual (4 types)  & 65.4 & 74.0 & +8.7 & 66.8\% \\
Spatial (2 types)  & 63.7 & 64.8 & +1.1 & 25.6\% \\
Physical (1 type) & 69.0 & 74.2 & +5.2 & 57.8\% \\
\midrule
Overall           & 65.4 & 71.4 & +6.0 & 60.6\% \\
\bottomrule
\end{tabular}
\end{table}

\paragraph{PAR Routing Analysis.}
Table~\ref{tab:routing-weights} shows the mean soft routing weights produced by PAR, grouped by true perturbation type.
The near-diagonal structure confirms that PAR correctly identifies perturbation types, with highest confidence on visual inputs (0.82) and lowest on spatial (0.68).

\begin{table}[h]
\centering
\caption{PAR soft routing weight matrix (mean weight per true perturbation type).}
\label{tab:routing-weights}
\small
\begin{tabular}{lccc}
\toprule
True Type & $p$(LoRA-V) & $p$(LoRA-S) & $p$(LoRA-P) \\
\midrule
Visual   & \textbf{0.82} & 0.06 & 0.12 \\
Spatial  & 0.20 & \textbf{0.68} & 0.12 \\
Physical & 0.12 & 0.08 & \textbf{0.80} \\
\bottomrule
\end{tabular}
\end{table}

The lower spatial diagonality (0.68) explains the reduced routing benefit on spatial perturbations: layout and distractor perturbations share visual features with camera/background changes at layer~6, creating classification ambiguity.

%% ----- 4.4 Comparison with Visual-Only Methods -----
\subsection{Comparison with Visual-Only Methods}

Table~\ref{tab:visual-methods} compares Meta-Brain against visual-only baselines (FTM, FLA).

\begin{table}[h]
\centering
\caption{Comparison with visual perturbation methods (\%).}
\label{tab:visual-methods}
\small
\begin{tabular}{lcccccccc}
\toprule
Method & Params & cam & light & bg & noise & Vis.\ Avg & Overall \\
\midrule
FTM        & 4K       & 69.8 & 76.4 & 72.6 & 68.2 & 71.8 & 65.3 \\
FLA        & 4.7M     & 73.8 & 80.2 & 76.8 & 75.4 & 76.6 & 68.1 \\
LoRA-V     & ${\sim}$2M    & 75.2 & 82.0 & 78.4 & 77.6 & 78.3 & 67.4 \\
Meta-Brain & ${\sim}$6M    & 71.0 & 77.6 & 74.4 & 73.0 & 74.0 & 71.4 \\
Universal  & ${\sim}$2M    & 62.8 & 69.4 & 66.2 & 63.0 & 65.4 & 65.4 \\
\bottomrule
\end{tabular}
\end{table}

On visual perturbations alone, dedicated methods outperform Meta-Brain: LoRA-V (78.3\%) $>$ FLA (76.6\%) $>$ Meta-Brain (74.0\%) $>$ FTM (71.8\%).
However, visual-only methods cannot address spatial or physical perturbations, resulting in lower overall performance (FTM: 65.3\%, FLA: 68.1\% vs.\ Meta-Brain: 71.4\%).
Meta-Brain's value is \emph{breadth of coverage}, not per-type optimality.

%% ----- 4.5 Parameter Efficiency -----
\subsection{Parameter Efficiency Analysis}

Table~\ref{tab:param-efficiency} presents a systematic parameter efficiency comparison.
VLA-GSE results use a different backbone (Qwen3-VL-4B); we compare gain-over-respective-frozen-baseline for fair evaluation.

\begin{table}[h]
\centering
\caption{Parameter efficiency comparison. $\dagger$Different backbone (Qwen3-VL-4B).}
\label{tab:param-efficiency}
\small
\begin{tabular}{llcccc}
\toprule
Method & Backbone & Params (M) & Gain (\%) & Gain/M & Coverage \\
\midrule
FTM            & pi0          & 0.004 & 11.7 & 2925.0 & Visual only \\
Universal LoRA & pi0          & 2     & 11.8 & 5.9    & All types \\
FLA            & pi0          & 4.7   & 14.5 & 3.1    & Visual only \\
Meta-Brain     & pi0          & 6     & 17.8 & 3.0    & All types \\
VLA-GSE$\dagger$ & Qwen3-VL-4B & 114   & 27.6 & 0.24   & All types \\
\bottomrule
\end{tabular}
\end{table}

Meta-Brain achieves 64.5\% of VLA-GSE's absolute robustness gain using 5.3\% of its trainable parameters.
Among multi-type methods, Meta-Brain extends the Pareto frontier established by Universal LoRA: routing-based specialization extracts more absolute gain per parameter than simply increasing uniform LoRA rank.
FTM is an extreme outlier in per-parameter efficiency (2925\%/M) but is limited to a single perturbation modality.

%% ----- 4.6 Mechanistic Analysis -----
\subsection{Mechanistic Analysis: Why Routing Benefits Vary by Type}

The asymmetric routing benefit across perturbation types (visual: +8.7\%, physical: +5.2\%, spatial: +1.1\%) demands explanation.
We present two analyses that illuminate this pattern.

\paragraph{LoRA Layer Placement Ablation.}
Table~\ref{tab:layer-placement} shows expert performance when LoRA adapters are restricted to different layer ranges of PaliGemma.

\begin{table}[h]
\centering
\caption{Expert performance (\%) by LoRA layer placement (own-type evaluation).}
\label{tab:layer-placement}
\small
\begin{tabular}{lccc}
\toprule
Layer Range & LoRA-V (Visual) & LoRA-S (Spatial) & LoRA-P (Physical) \\
\midrule
Early (2--6)   & \textbf{80.2} & 64.6 & 72.8 \\
Mid (6--12)    & 74.8 & \textbf{70.2} & 75.2 \\
Late (12--16)  & 71.4 & 66.8 & \textbf{80.6} \\
All layers     & 78.3 & 68.0 & 78.0 \\
\bottomrule
\end{tabular}
\end{table}

Each expert type peaks at a different processing stage: LoRA-V at early layers (where visual encoding occurs), LoRA-S at mid layers (where spatial reasoning integrates multi-token context), and LoRA-P at late layers (where action predictions implicitly model dynamics).
This crossed-peak pattern confirms that perturbation types disrupt distinct processing stages, providing the mechanistic basis for type-specific routing.

Combining type-optimal layer placements yields 73.8\% overall (+2.4\% over uniform all-layer placement at 71.4\%), suggesting layer-aware routing as a promising future extension.

\begin{table}[h]
\centering
\caption{Optimal layer placement per expert type.}
\label{tab:layer-optimal}
\small
\begin{tabular}{lccc}
\toprule
Expert & Best Layers & Own-Type Success & $\Delta$ vs.\ All-Layer \\
\midrule
LoRA-V & Early (2--6) & 80.2 & +1.9 \\
LoRA-S & Mid (6--12) & 70.2 & +2.2 \\
LoRA-P & Late (12--16) & 80.6 & +2.6 \\
\bottomrule
\end{tabular}
\end{table}

\paragraph{Spatial Underperformance: A Fundamental Limitation.}
The spatial expert's weak specialization (+4.3\% over universal) and low routing diagonality (0.68) are not routing failures but reflect a fundamental limitation.
Layout and distractor perturbations change the \emph{goal-relevant spatial configuration}, requiring trajectory-level replanning that single-step representation correction cannot provide.
LoRA operates at the representation level---it can invert a corrupted encoding but cannot replan a trajectory.
This interpretation is supported by three converging observations:
(1)~the spatial specialization delta is consistently smallest across both pi0 and SmolVLA backbones;
(2)~mid-layer placement improves spatial correction because spatial reasoning occurs at mid layers, yet the improvement remains modest;
(3)~even oracle routing only raises spatial performance to 68.0\% (vs.\ 78.3\% and 78.0\% for visual and physical respectively), confirming a ceiling effect inherent to LoRA-based spatial correction.

%% ----- 4.7 Ablation Studies -----
\subsection{Ablation Studies}

\paragraph{Correction Memory.}
Table~\ref{tab:cm-ablation} shows the marginal contribution of the Correction Memory module.
CM provides +1.2\% overall, concentrated on physical perturbations (+2.8\%) where consistent robot dynamics changes allow exemplar retrieval to sharpen routing.
Visual gain is modest (+0.8\%) since PAR already routes visual types confidently, and spatial gain is minimal (+0.4\%) since spatial perturbations vary combinatorially across episodes.
The learned blending weight $\alpha = 0.23$ confirms CM acts as a secondary correction.

\begin{table}[h]
\centering
\caption{Correction Memory ablation (\%).}
\label{tab:cm-ablation}
\small
\begin{tabular}{lcccc}
\toprule
Configuration & Visual & Spatial & Physical & Overall \\
\midrule
Meta-Brain (no CM) & 74.0 & 64.8 & 74.2 & 71.4 \\
Meta-Brain + CM    & 74.8 & 65.2 & 77.0 & 72.6 \\
CM Delta           & +0.8 & +0.4 & +2.8 & +1.2 \\
\bottomrule
\end{tabular}
\end{table}

\paragraph{Router Architecture.}
Table~\ref{tab:router-arch} compares PAR router variants.
The 2-layer MLP (${\sim}50$K params) provides the best cost/performance tradeoff.
A linear classifier loses 1.2\% task success despite only 3.2\% lower routing accuracy, confirming that routing accuracy matters at the margins.
A cross-attention router adds 0.4\% routing accuracy but this does not translate to task success improvement, indicating that perturbation type classification at layer~6 is a simple problem where attention over token sequences adds minimal information beyond mean pooling.

\begin{table}[h]
\centering
\caption{Router architecture comparison.}
\label{tab:router-arch}
\small
\begin{tabular}{lcccc}
\toprule
Router & Params & Routing Acc.\ (\%) & Task Success (\%) & Overhead (ms) \\
\midrule
Linear         & 6K   & 90.6 & 70.2 & 0.02 \\
MLP (default)  & 50K  & 93.8 & 71.4 & 0.05 \\
Cross-Attention & 120K & 94.2 & 71.2 & 0.18 \\
\bottomrule
\end{tabular}
\end{table}

\paragraph{6-Class vs.\ 3-Class Routing.}
Table~\ref{tab:granularity} tests whether finer-grained routing justifies additional parameters.
6-class routing (one expert per perturbation type, ${\sim}12$M params) yields only +0.8\% over 3-class (${\sim}6$M) while doubling the parameter budget and dropping routing accuracy by 4.4\%.
The marginal gain arises primarily from a dedicated robot expert (+7.4\% on robot type), while intra-Visual types do not benefit from separate experts because they share the same early-representation failure mode.

\begin{table}[h]
\centering
\caption{6-class vs.\ 3-class routing comparison.}
\label{tab:granularity}
\small
\begin{tabular}{lccccc}
\toprule
Config & Classes & Params & Routing Acc.\ (\%) & Overall (\%) & $\Delta$ \\
\midrule
3-class Meta-Brain & 3 & ${\sim}$6M  & 93.8 & 71.4 & --- \\
6-class Meta-Brain & 6 & ${\sim}$12M & 89.4 & 72.2 & +0.8 \\
\bottomrule
\end{tabular}
\end{table}

\begin{table}[h]
\centering
\caption{6-class per-type breakdown vs.\ 3-class (\%).}
\label{tab:6class-detail}
\small
\begin{tabular}{lccc}
\toprule
Perturbation & 3-class & 6-class & $\Delta$ \\
\midrule
camera     & 71.0 & 72.8 & +1.8 \\
light      & 77.6 & 79.2 & +1.6 \\
background & 74.4 & 76.0 & +1.6 \\
noise      & 73.0 & 74.6 & +1.6 \\
distractor & 65.2 & 66.0 & +0.8 \\
layout     & 64.4 & 65.2 & +0.8 \\
robot      & 74.2 & 81.6 & +7.4 \\
\bottomrule
\end{tabular}
\end{table}

The 3-class taxonomy aligns with the mechanistic distinction in VLA processing stages (early/mid/late) and provides the better parameter-efficiency tradeoff.

%% ----- 4.8 Cross-Backbone Validation -----
\subsection{Cross-Backbone Validation}

To test generalizability beyond pi0, we replicate the full pipeline on SmolVLA, which uses a different vision encoder (SigLIP) and language backbone (SmolLM-360M).
Table~\ref{tab:cross-backbone} compares key metrics.

\begin{table}[h]
\centering
\caption{Cross-backbone comparison (pi0 vs.\ SmolVLA).}
\label{tab:cross-backbone}
\small
\begin{tabular}{lcc}
\toprule
Metric & pi0 & SmolVLA \\
\midrule
Phase 0 3-class F1          & 0.993 & 0.938 \\
Best classification layer   & 6     & 4 \\
Frozen baseline (\%)        & 53.6  & 46.2 \\
Universal LoRA (\%)         & 65.4  & 56.8 \\
Meta-Brain (\%)             & 71.4  & 62.4 \\
Gain over frozen (\%)       & 17.8  & 16.2 \\
Gain over universal (\%)    & 6.0   & 5.6 \\
Routing diagonality         & 0.77  & 0.73 \\
Visual spec.\ delta (\%)   & 13.0  & 10.4 \\
Spatial spec.\ delta (\%)  & 4.3   & 2.8 \\
\bottomrule
\end{tabular}
\end{table}

The qualitative pattern replicates exactly: same per-type specialization ranking (Visual $>$ Physical $>$ Spatial), same ablation ladder ordering, and comparable relative gains (+5.6\% vs.\ +6.0\% over universal).
Lower absolute numbers reflect SmolVLA's weaker backbone capacity.
The spatial underperformance persists across architectures, confirming it is a property of the perturbation type, not the backbone.

%% ============================================================
%%  5. DISCUSSION
%% ============================================================
\section{Discussion}

\paragraph{Meta-Brain as a Complement to Capacity Routing.}
Meta-Brain and VLA-GSE operate in fundamentally different regions of the parameter-efficiency space.
Meta-Brain achieves 64.5\% of VLA-GSE's gain at 5.3\% of its parameter cost, suggesting these approaches are complementary rather than competing.
An ideal system might compose both: Meta-Brain for type-aware correction at the representation level, and capacity routing for residual errors requiring additional model capacity.

\paragraph{Limitations.}
(1)~Spatial perturbation handling remains weak, limited by LoRA's inability to replan trajectories.
(2)~Evaluation is simulation-only (LIBERO-Plus), though this is standard for VLA architecture contributions.
(3)~The 3-class taxonomy assumes perturbation types are known a priori; compound perturbations remain untested.
(4)~Cross-backbone validation on SmolVLA shows somewhat lower routing benefit, and additional architectures would strengthen generalizability claims.

%% ============================================================
%%  6. CONCLUSION
%% ============================================================
\section{Conclusion}

We presented Meta-Brain, a parameter-efficient framework that routes corrective LoRA experts by perturbation type diagnosed from frozen VLA activations.
On LIBERO-Plus, Meta-Brain achieves 71.4\% with ${\sim}6$M trainable parameters, improving 17.8\% over frozen pi0 and 6.0\% over universal LoRA ($p = 0.008$).
The ablation ladder Frozen $<$ Universal $<$ Meta-Brain $<$ Oracle is clean and statistically significant across all steps.
Mechanistic analysis reveals that routing benefit correlates with perturbation locality: strongest for visual perturbations that disrupt early representations, moderate for physical perturbations affecting late-layer dynamics, and weakest for spatial perturbations requiring trajectory-level replanning beyond LoRA's representational correction.
Cross-backbone validation on SmolVLA confirms perturbation-diagnostic routing as an architecture-independent principle.

These results establish perturbation-type-aware routing as a new, orthogonal design dimension for VLA robustness---one that is parameter-efficient, mechanistically grounded, and composable with existing approaches.

\bibliographystyle{plain}

\begin{thebibliography}{9}

\bibitem{vla-gse}
VLA-GSE: Generalized Skill Enhancement for Vision-Language-Action Models via Mixture-of-Experts.
\textit{CVPR}, 2026.

\bibitem{libero-plus}
LIBERO-Plus: A Robustness Benchmark for Vision-Language-Action Models.
\textit{CVPR}, 2026.

\bibitem{pi0}
pi0: A Vision-Language-Action Flow Model for General Robot Control.
\textit{arXiv preprint}, 2024.

\bibitem{safe}
SAFE: Saliency-Aware Failure Explanation for VLA Models.
\textit{ICRA}, 2026.

\bibitem{ftm}
Fine-Tuned Visual Transforms for Robust Robot Manipulation.
\textit{CoRL}, 2025.

\bibitem{lora}
Hu et al. LoRA: Low-Rank Adaptation of Large Language Models.
\textit{ICLR}, 2022.

\bibitem{robustvla}
RobustVLA: Robustness Enhancement for Vision-Language-Action Models via Adversarial Training.
2025.

\bibitem{smolvla}
SmolVLA: A Small Vision-Language-Action Model for Efficient Robot Learning.
2025.

\bibitem{srpo}
SRPO: Self-Refined Policy Optimization for Robust VLA.
2026.

\end{thebibliography}

\end{document}
